% explainer-source-hash: 3545e9e2112977ac2ee390e08fc010506ae5e423e66d118f3678abcc3a4f799e
% explainer-source-commit: a3b383a90a485c793e6efa6e9587845e06db8a96
% explainer-generated: 2026-07-16
% Regenerate with /explain-all (see WIRING.md). Do not edit this stamp by hand:
% it is how the toolchain knows whether this document still matches the code.
\documentclass[11pt,a4paper]{article}
\input{preamble}

\title{\vspace{-1.5cm}\textbf{Custom Web Scraper}\\[0.3cm]
\large The Brief: what it is, how it is built, where it is weak}
\author{Portfolio explainer, built from the source}
\date{}

\begin{document}
\maketitle
\thispagestyle{empty}

\section*{In one paragraph}

\begin{keyidea}{The project}
A desktop application that opens a real Chrome browser, loads a web address you
type, and then lets you interrogate that live page with any of seven Selenium
locator strategies, showing every match in a table and outlining them in the
browser itself. It exists because small scraping investigations do not deserve a
throwaway Python script each: load the page once, then fire locator after locator
at it from a graphical interface. Roughly 550 lines of Python in one file, with
exactly one third-party dependency.
\end{keyidea}

\section{Why it is built the way it is}

Two decisions define the project, and both are defensible.

\textbf{It drives a real browser instead of downloading HTML.} The simple approach,
fetching a page over HTTP and parsing the text, sees only the initial response.
Most modern sites build their content afterwards with JavaScript running inside
the browser, so that approach returns an empty shell. Selenium drives actual
Chrome, so the page is fully rendered before anything is read. The price is a
Chrome install and a window on your screen, and the project accepts it knowingly.

\textbf{It is a GUI, not a library.} The target user is someone poking at an
unfamiliar page for five minutes. The interface encodes that workflow directly:
panel ``1. Load a page'', panel ``2. Find elements''. The order of operations is
the layout.

\section{Architecture}

Three layers live in one file, separated by convention rather than by folders.

\begin{figure}[H]
\centering
\resizebox{0.95\textwidth}{!}{
\begin{tikzpicture}[node distance=8mm]

\node[box, minimum width=36mm] (ui) {\textbf{Tkinter layer}\\[2pt]
  \scriptsize builds the window,\\
  \scriptsize routes button clicks,\\
  \scriptsize renders the results table};

\node[box, minimum width=36mm, right=30mm of ui] (sel) {\textbf{Selenium layer}\\[2pt]
  \scriptsize owns the browser,\\
  \scriptsize navigates and waits,\\
  \scriptsize runs the seven lookups};

\node[box, minimum width=36mm, below=15mm of ui] (pure) {\textbf{Pure formatting layer}\\[2pt]
  \scriptsize elements $\rightarrow$ plain tuples,\\
  \scriptsize no driver, no globals,\\
  \scriptsize testable in milliseconds};

\node[extbox, minimum width=28mm, right=30mm of sel] (chrome) {\textbf{Real Chrome}};
\node[extbox, minimum width=28mm, below=15mm of chrome] (web) {\textbf{The website}};

\node[store, below=15mm of sel, minimum width=20mm] (state) {\scriptsize module globals\\\scriptsize \code{driver}\\\scriptsize \code{is\_request\_loaded}\\\scriptsize \code{last\_elements}};

\draw[flow] (ui) -- node[lbl, above]{callbacks} (sel);
\draw[flow] (sel) -- node[lbl, above]{WebDriver} (chrome);
\draw[flow] (chrome) -- node[lbl, right]{HTTP} (web);
\draw[flow] (sel) -- node[lbl, right]{elements} (state);
\draw[dflow] (state) -- (pure);
\draw[dflow] (pure) -- node[lbl, left]{rows} (ui);

\begin{scope}[on background layer]
\node[draw=linegrey, fill=softgrey, rounded corners, fit=(ui)(pure)(sel)(state),
      inner sep=7mm, label={[font=\small\itshape, text=inkblue]above:main.py}] {};
\end{scope}
\end{tikzpicture}}
\caption{Three layers, three shared globals, two things outside the process.}
\end{figure}

The \textbf{pure formatting layer} is the part worth pointing at. It converts
Selenium's live element handles into flat display tuples, and it does so without
touching the driver or any global: it only calls \code{.tag\_name}, \code{.text}
and \code{.get\_attribute()} on whatever object it is handed. Any object with those
three members will do, so it can be tested with a five-line fake class, no browser
and no network. That is a real design decision, and it is the strongest thing in
the codebase.

The rest of the file cannot make that claim, and the reason is the box in the
middle of the diagram: \code{driver}, \code{is\_request\_loaded} and
\code{last\_elements} are module-level globals that six functions reach into
directly.

\section{How one lookup flows}

\begin{figure}[H]
\centering
\resizebox{\textwidth}{!}{
\begin{tikzpicture}[node distance=7mm]
\node[box] (url) {Type a URL,\\click Request};
\node[box, right=12mm of url] (get) {\code{driver.get(url)}\\wait for html, body};
\node[box, right=12mm of get] (pick) {Pick a strategy\\+ type a value};
\node[dec, right=13mm of pick] (d) {locator\\matches?};
\node[box, below=12mm of d] (find) {\code{find\_elements()}\\returns all matches};
\node[box, left=13mm of find] (rows) {format to tuples};
\node[box, left=13mm of rows] (win) {new window\\with a table};
\node[box, fill=white, draw=warnred, text=warnred, right=13mm of d] (err) {status bar:\\``enter a valid X''};

\draw[flow] (url) -- (get);
\draw[flow] (get) -- (pick);
\draw[flow] (pick) -- (d);
\draw[flow] (d) -- node[lbl, right]{yes, within 10s} (find);
\draw[flow] (find) -- (rows);
\draw[flow] (rows) -- (win);
\draw[flow, draw=warnred] (d) -- node[lbl, above]{timeout} (err);
\end{tikzpicture}}
\caption{The path of a single lookup. The ``timeout'' branch also fires when a
perfectly valid locator simply matches nothing, which is the source of the
mislabelled-error defect in Section~\ref{sec:weak}.}
\end{figure}

The seven strategies are class, ID, name, tag name, link text, partial link text
and XPath, held in a \code{LOCATOR\_DISPATCH} dictionary keyed by the dropdown's
label. The dropdown's initial value, \code{"Select"}, is deliberately absent from
that dictionary, so a missing key doubles as the ``you have not picked one yet''
check with no second list to keep in sync.

Each lookup waits before it finds. A ten-second
\code{WebDriverWait} on \code{presence\_of\_element\_located} blocks until at
least one match exists, and then \code{find\_elements} collects all of them.
This is the correct pattern for a page that may still be rendering, and the
project uses it consistently rather than sleeping and hoping.

\section{Features of substance}

\begin{itemize}
\item \textbf{Highlight in Browser.} After a lookup, this injects a few lines of
  JavaScript into the live page via \code{execute\_script}, outlining each matched
  element in pink and scrolling it into view. It closes the loop between ``my
  locator matched 12 things'' and ``are they the 12 I meant?'', which is the
  entire question a locator-testing tool exists to answer. The best feature here.
\item \textbf{Status bar.} A single line, recoloured by severity from a three-key
  dictionary, reporting every step. A \code{root.update\_idletasks()} call forces
  the ``Loading page\ldots'' repaint before the browser blocks the event loop, so
  the window does not simply appear to freeze.
\item \textbf{Session history.} Both text fields are comboboxes backed by one
  shared 10-entry most-recently-used list. In memory only; it dies with the
  process, which the README states plainly.
\item \textbf{Contextual instructions.} A label under the dropdown rewrites itself
  per strategy through a \code{StringVar} trace callback.
\item \textbf{Portable driver resolution.} \code{CHROMEDRIVER\_PATH} points at a
  specific binary; unset, Selenium Manager fetches the right one automatically.
  This replaced a driver path hardcoded to a Windows folder on the original
  author's machine, which meant the project ran on exactly one computer on
  earth. It is the best ``portability problem I solved'' story here.
\end{itemize}

\section{Where it is weak}
\label{sec:weak}

\begin{honest}{The README claims a test suite the repository does not contain}
The README states that a 7-case unit suite and a full window test pass, and
refers the reader to a ``Test plan'' section for the command and output. That
section does not exist, and no test file has ever been committed: \code{git
ls-files} lists eight files, none a test, and a search of every commit for every
file ever added confirms it was never there.

The claims themselves appear to be true. Equivalent tests written from scratch for
this document reproduced every stated behaviour. So the tests were evidently
written and run but never committed, and the dangling cross-reference is what
exposes it. This is the project's most urgent problem, and it is not in the code:
an interviewer who reads ``all 7 cases pass; see Test plan below'' will go looking
and find nothing. Commit the tests or drop the claim.
\end{honest}

\begin{honest}{Two further README claims the code contradicts}
\textbf{The results table is not sortable}, though the opening paragraph says it
is. \code{ttk.Treeview} sorts only when a command is bound to each column heading,
and none is bound.

\textbf{\code{temp.py} and \code{temp1.py} are not ``still visible in the source
history''.} The BeautifulSoup draft they refer to may have existed locally, but
this repository has four commits and has never contained a file by either name.
\end{honest}

\begin{honest}{Real defects in the code}
\begin{itemize}
\item \textbf{``No matches'' is reported as ``please enter a valid class name''.}
  Because the wait must succeed before \code{find\_elements} runs, a locator that
  matches nothing always times out first. The zero-match handling that follows,
  the count of 0, the disabled button, the ``No elements matched'' label, is
  therefore unreachable, and a valid locator gets blamed on a typo the user did
  not make. Confirmed by forcing the branch with a stubbed wait: the code is
  correct, just unreachable.
\item \textbf{\code{is\_request\_loaded} is set before the load is confirmed}, so
  a page-load timeout leaves the app believing a page is ready when it is not.
\item \textbf{Only \code{InvalidArgumentException} is caught around
  \code{driver.get}}, so a dead host or crashed browser escapes as a raw traceback
  while the GUI sits on ``Loading page\ldots''. For a tool whose selling point is
  never failing silently, that is the gap.
\item \textbf{Stale highlight state.} \code{last\_elements} and the highlight
  button are never reset on navigation or after a failed lookup, so the button
  stays enabled pointing at a previous page's elements.
\item \textbf{Every lookup opens a new window} and none is ever closed, so the
  core workflow of trying locator after locator is also what buries the screen.
\end{itemize}
\end{honest}

\begin{honest}{The globals, and the seven near-identical functions}
The seven \code{get\_element\_by\_*} functions differ only in one \code{By}
constant and one error string; the guard, the wait, the exception handler and the
return shape are copied seven times. \code{LOCATOR\_DISPATCH} removed the
\code{if}/\code{elif} chain in the caller but left the duplication in the callees.

All seven read \code{driver} from module scope rather than receiving it, so none
can be tested without patching the module, and the app structurally cannot hold
two sessions. \code{driver} is not even given a module-level default; it exists
only once the \code{\_\_main\_\_} block assigns it, so importing the module and
calling anything raises \code{NameError} until a fake is patched in.

The README says all of this about itself, unprompted, which counts for something.
The fair framing is not ``this is wrong'' but ``this traded testability for
brevity in a 550-line single-window tool, knowingly''. The fix is a small class
holding \code{driver} and \code{is\_request\_loaded} as instance attributes.
\end{honest}

\section{What it gets right}

\begin{itemize}
\item The pure formatting layer, and specifically the choice to wrap \emph{each
  field access} in a small exception-swallowing helper rather than the whole row,
  so an element going stale mid-read still yields a partial row. That granularity
  shows the stale-element problem was understood, not just caught.
\item Truncation that is actually correct: 200 characters in gives exactly 80 out,
  ellipsis included in the budget rather than pushed past it. Verified.
\item A dependency range chosen for a reason: \code{selenium>=4.15,<5}, floored
  above the 4.6 release that introduced Selenium Manager, ceilinged below the next
  major version's licence to break the API.
\item Toolkit fluency in small places: \code{update\_idletasks} before a blocking
  call, \code{rowconfigure} weights so the table absorbs resize, the \code{clam}
  theme wrapped in a \code{try}/\code{except} because it is not on every platform.
\item One genuinely subtle guard, \code{if "instructions" not in globals():
  return}, which looks like dead paranoia and is not: constructing the dropdown
  writes its variable and fires the trace callback immediately, tens of lines
  before the label it updates is created. Without it, startup raises
  \code{NameError} every time.
\end{itemize}

\section{Verified for this document}

\begin{itemize}
\item \code{py\_compile} clean on Python 3.12.3 with selenium 4.45.0.
\item The formatting functions driven with fake duck-typed elements: truncation,
  whitespace collapsing, attribute filtering, 1-based numbering, and a fully stale
  element yielding \code{(3, '?', '', '')}. All as described.
\item The entire window built and driven on a real display against a fake driver:
  load, lookup, results table, and highlight all produced the expected status text
  and the expected 2 \code{execute\_script} calls.
\item Repository history searched: no test file, no \code{temp.py}, ever.
\end{itemize}

\textbf{Not verified:} a live Chrome session against a real site, and the
highlight JavaScript against a real DOM. Both need Chrome plus network access for
Selenium Manager. The Selenium API used is standard 4.x, but that claim rests on
inspection, not execution.

\section*{Glossary}

\begin{description}[leftmargin=0pt, style=nextline, itemsep=3pt]
\item[DOM] The browser's live in-memory tree of a page, built from its HTML. What
  Selenium actually queries.
\item[Element] One node of that tree: a link, a paragraph, an image.
\item[Locator] A rule for finding elements, in two halves: a strategy (by class?
  by ID?) and a value (which class?). The split the whole interface is built on.
\item[Selenium / WebDriver] A library, and the standard protocol it speaks, for
  remote-controlling a real browser from code.
\item[Selenium Manager] Bundled since selenium 4.6; detects your Chrome and
  downloads a matching driver, removing the version-matching chore by hand.
\item[Tkinter / ttk] Python's built-in GUI toolkit, and its themed widget set.
\item[Event loop] The endless wait-for-input, call-the-callback cycle a GUI runs
  in. \code{root.mainloop()} enters it and does not return until the window closes.
\item[Callback] A function registered to run when something happens, such as a
  button being clicked.
\item[Duck typing] Python's rule that only an object's methods matter, not its
  type. What makes the formatting layer testable with fakes.
\item[Stale element reference] The error raised when you read an element handle
  after the page has re-rendered underneath it. Routine on modern sites.
\item[XPath] A query language for walking a document tree, and the escape hatch
  that expresses what the other six strategies cannot.
\item[Headless] Running the browser with no visible window. Enabled here by
  setting \code{SCRAPER\_HEADLESS=1}.
\end{description}

\end{document}
